%% CPP 2027 Submission — Formalizing Operator Semantics for Quantized Recurrent Inference
%% Compile: pdflatex + bibtex (acmart)
\documentclass[sigplan,10pt,anonymous,review]{acmart}
\settopmatter{printfolios=true,printacmref=false}

\AtBeginDocument{%
  \providecommand\BibTeX{{\normalfont B\scshape ib\TeX}}}

\settopmatter{printacmref=false}
\renewcommand\footnotetextcopyrightpermission[1]{}
\pagestyle{plain}

\usepackage[utf8]{inputenc}
\usepackage[T1]{fontenc}
\usepackage{amsmath}
\ifdefined\XeTeXinterchartokenstate
  % XeLaTeX: unicode-math (loaded by acmart) supplies the amssymb glyphs
\else
  \usepackage{amssymb}
\fi
\usepackage{booktabs}
\usepackage{microtype}

\begin{document}

\title{Formalizing Operator Semantics for Quantized Recurrent Inference:\\
basepoint-relative operators, symplectic projection,\\
and Q2\_0 symbol quantization (C026--C038, all PROVED, no \texttt{sorry})}

\author{Anonymous}
\affiliation{%
  \institution{Independent}
  \country{}}


\date{2026-08-20}

\begin{abstract}
We report a Lean 4 + mathlib formalization of operator semantics for quantized
recurrent inference (claims C026--C038; all PROVED; novelty: KNOWN except C037
(unassessed); no \texttt{sorry}; full \texttt{lake build} passes). The formalization
is intuition-guided: every operator carries its own basepoint (the additive
operator is anchored at $0$, the multiplicative operator at $1$, so ``$2\times2$''
is not the $2$ of the real axis); the multiplicative basepoints $\pm1$ are shown
to be the projection extremes of a higher-dimensional structure (the unit circle),
which explains the four-phase discrete group $\{1,-1,i,-i\}$ under the symplectic
rotation $J^2=-I$; the Q2\_0 quantization format of GGUF weights (grid
$\{-1,0,1,2\}$, 2 bits per element) is formalized as a \emph{symbol encoding} of
the four phases by two sign axes (real $\times$ imaginary), not as an algebraic
homomorphism; and the SSM state matrix is shown to carry a complex structure
(each $2\times2$ block $[a\ b;-b\ a] = aI + bJ$ is one complex number $a+bi$
stored with 4 parameters --- 50\% square redundancy), with a lossless complex
compactification. We bridge the formalization to an empirical inference backend:
an assembly kernel chain that implements the same semantics achieves
bit-exact alignment with a float reference on the input chain and preserves
locked-domain decisions end-to-end on real Bonsai-27B weights.
\end{abstract}

\keywords{Lean, formalization, quantization, SSM, basepoint, symplectic, operator semantics}

\maketitle

\section{Introduction}

Quantized recurrent inference --- state-space models (SSMs) with 2-bit GGUF
weights --- is a production reality, yet the \emph{operator semantics} of the
quantization grid is rarely formalized: what does the grid $\{-1,0,1,2\}$
\emph{mean}? This paper formalizes an answer: the grid is a symbol encoding of
the four-phase structure $\{1,-1,i,-i\}$ of the complex plane, and the SSM state
matrix carries a complex structure that makes half of its storage redundant.

The formalization chain (C026--C038) starts from a basepoint-relative view of
operators (C026): every operator has its own basepoint, so iteration semantics
is basepoint-dependent. The multiplicative basepoints $\pm1$ are then explained
as projection extremes of the unit circle (C029), yielding the four-phase
discrete group under the symplectic rotation $J^2=-I$ (C030, C031). The dual
projection (C033) reads wave--particle duality as the $\mathrm{Re}/\mathrm{Im}$
projections of one complex amplitude; matrix operators (C035) and
trigonometric direction (C036) follow the same axis semantics. Two claims bridge
to engineering: C037 formalizes the Q2\_0 quantization grid as a sign encoding
of the four phases, and C038 formalizes the complex structure of SSM states
with its 50\% redundancy and lossless compactification.

All claims are machine-checked in Lean 4 (mathlib v4.32.2), with no
\texttt{sorry}. The artifact includes a fully buildable Lean project and an
empirical bridge: an assembly kernel chain implementing the same semantics
(weight locking per block, window sliding with pre-locking, integer-MAC
convolution, L2 compactification, interlock update) achieves bit-exact
alignment with a numpy pat0 reference on the input chain (relative error
$<10^{-4}$) and preserves locked-domain decisions end-to-end on real
Bonsai-27B Q2\_0 weights.

\paragraph{Contributions.}
\begin{enumerate}
\item A complete Lean formalization (C026--C038, no \texttt{sorry}) of
basepoint-relative operator semantics (\S\S3), the four-phase discrete group and
symplectic projection (\S\S3.2--3.3), Q2\_0 symbol quantization (\S4), and SSM
state complex compactification (\S5).
\item The Q2\_0 four-point grid as a two-sign-axis encoding of the four phases:
the grid is a \emph{symbol encoding}, not a homomorphism (2$\cdot$2 = 4 leaves
the grid).
\item A lossless complex compactification of SSM states (50\% square redundancy)
with an empirical deviation measurement ($\approx 3\times10^{-5}$).
\item An empirical bridge: an assembly inference chain whose input chain is
bit-exact against a float reference, with locked-domain decision preservation
on real weights.
\end{enumerate}

\section{Background: basepoints, heaps, and iteration}

\subsection{Basepoint-relative algebra}

The heap/torsor theory of basepoint-relative algebra was formalized previously
(C001--C010): choosing a basepoint $e$ in a heap recovers a group (C004); a
semiconjugacy preserves iteration and basepoint orbits (C001); the basepoint-relative
step translation $\tau_{e,a}(x) = [x,e,a]$ transports under basepoint change only
if the endpoint is transported (C006); the invariant is the relative displacement
$(e,a)$, not the point $a$ (C007). Pure-heap generation laws are
transport-covariant (C010/R011): structure bifurcation requires breaking that
covariance.

\subsection{From heaps to operators}

The operator layer (C026) applies the same basepoint discipline to operators:
each operator is anchored at its own basepoint, so iteration is
basepoint-relative. The multiplicative basepoint $1$ is \emph{not} the additive
origin $0$; a basepoint move conjugates the additive operator to itself but
turns a multiplicative operator into an affine map. This is the algebraic core
that the rest of the paper builds on.

\section{Operator semantics (C026--C036)}

\subsection{Every operator has its own basepoint (C026, C028)}

\begin{theorem}[C026]
On any commutative ring $R$: with $s(x) = x+1$ and $m_k(x) = k\cdot x$,
$m_2(1) = 2 = s(s(0))$, but $m_2(m_2(1)) = 4 \ne 2 = s(s(0))$ in $\mathbb Z$.
The basepoint move $T(x) = x-1$ maps $1 \mapsto 0$; $T\circ s \circ T^{-1} = s$;
$T\circ m_2 \circ T^{-1}(y) = 2y + 1$ (multiplication becomes affine); the
conjugated operator's fixed point is $-1 \ne 0$.
\end{theorem}

The full operator family (C028) anchors each operator at its basepoint:
addStep at $0$ (arithmetic chain $0,1,2,\dots$), mulStep at $1$ (power chain
$1,k,k^2,\dots$), negOp at $0$ (involution), recipOp at $\pm1$ (involution),
sqrtOp at $1$ (contraction). The Q2\_0 grid $\{-1,0,1,2\}$ is read as the
basepoint cluster $\{\pm1, 0, 2\}$: every quantized weight lands on the
nearest basepoint or basepoint iteration.

\subsection{$\pm1$ = projection of a higher-dimensional structure (C029)}

\begin{theorem}[C029]
On the unit circle $S^1 = \{z : \mathbb C \mid \|z\|=1\}$, the real projection
$\mathrm{Re}$ has image $[-1,1]$; the points attaining the extremes are exactly
$z = \pm1$. The four phases $\{1,-1,i,-i\}$ project to $\{\pm1\}$ on the real
axis and $\{\pm i\}$ on the imaginary axis.
\end{theorem}

Two basepoint positions $\pm1$ cannot arise internally on a one-dimensional
axis; they are the projection extremes of the unit circle. This is the
projection reading of the multiplicative basepoints, and it explains why the
Q2\_0 grid contains $\pm1$.

\subsection{Symplectic projection and the four-phase group (C030, C031)}

\begin{theorem}[C030]
With $J(z) = i\cdot z$ ($J^2 = -I$), the four phases cycle
$1 \to i \to -1 \to -i \to 1$. Their vector sum is $0$; for every direction $u$,
$\sum_{p \in \{1,i,-1,-i\}} \mathrm{Re}(p\cdot\overline{u}) = 0$
(projection cancellation in every direction).
\end{theorem}

\begin{theorem}[C031]
The four phases form the $J$-orbit $C_4$. The reciprocal map equals
conjugation on the unit circle ($\|z\|=1 \Rightarrow z^{-1} = \overline{z}$)
and anticommutes with $J$ ($\overline{Jz} = -J\overline{z}$), extending each
phase pair to the dihedral group $D_4$ (4 rotations $\times$ 2 reflections).
The basepoint cluster is $\{0, \pm1, \pm i\}$ (5 positions).
\end{theorem}

\subsection{Dual projection (C033)}

\begin{theorem}[C033]
For $z \in \mathbb C$: $z = \mathrm{Re}\,z + i\,\mathrm{Im}\,z$ (unique
decomposition); $J$ swaps the projections ($\mathrm{Re}(Jz) = -\mathrm{Im}\,z$,
$\mathrm{Im}(Jz) = \mathrm{Re}\,z$); the two projection axes are orthogonal.
The four phases are the pure states: $1$ = pure particle, $i$ = pure wave,
$-1$ = anti-particle, $-i$ = anti-wave.
\end{theorem}

\subsection{Matrix and trigonometric direction (C035, C036)}

Matrix operations are read axis-semantically (C035): the dot product is axis
contraction (basepoint $0$); the cross product is interlock generating the
orthogonal complement axis ($\hat e_x \times \hat e_y = \hat e_z$, self-cross $=0$);
the outer product is the interlock matrix (the SSM state update $k\otimes d$);
matrix multiplication is shared-axis contraction whose correctness depends on
axis alignment (a transposed layout is an axis misalignment).

Trigonometric functions are a direction problem (C036):
$\cos(-\theta) = \cos\theta$ (even: direction lost) while
$\sin(-\theta) = -\sin\theta$ (odd: direction kept); their basepoints differ
($\cos 0 = 1$ vs $\sin 0 = 0$); the $90^\circ$ grid takes
$(\cos,\sin) \in \{(1,0),(0,1),(-1,0),(0,-1)\}$, isomorphic to the symplectic
rotation ($\mathrm{rotJ}^2 = -I$, $\mathrm{rotJ}^4 = I$).

\section{Q2\_0 symbol quantization (C037)}

The central bridging claim. GGUF Q2\_0 quantization stores 2 bits per element
with grid $\{-1, 0, 1, 2\}$; pat0's native numeric grid has 3 values
$\{-1,0,1\}$. The count difference is a projection direction:

\begin{theorem}[C037]
Let $\varphi : \mathbb C \to \mathbb R$ be the sign quantization
$\varphi(1)=1$, $\varphi(-1)=-1$, $\varphi(i)=2$, $\varphi(-i)=0$. Then:
(a) the images of the four phases are exactly the Q2\_0 grid
$\{-1,0,1,2\}$; (b) $1 + i + (-1) + (-i) = 0$ (projection cancellation);
(c) $\varphi$ is \emph{not} a homomorphism: $2\cdot 2 = 4 \notin \{-1,0,1,2\}$
(the grid is not multiplicatively closed) --- $\varphi$ preserves sign
directions, not multiplication; (d) the pat three-grid $\{-1,0,1\}$ is the
real-axis projection of the four phases ($\mathrm{Re}(\pm i) = 0$: imaginary
sign dropped).
\end{theorem}

The Q2\_0 grid is thus a two-sign-axis encoding (real sign $\times$ imaginary
sign): 2 bits = 2 sign bits, preserving the imaginary sign that the pat
three-grid drops. Positive imaginary ($i$) maps to $+2$ (basepoint iteration
$1+1 = 2\cdot 1$, C028); negative imaginary ($-i$) folds into $0$ (the fold
center). This gives the quantization format an operator-semantic meaning:
each weight is quantized to the nearest basepoint or basepoint iteration.

\section{State-matrix complex structure (C038)}

\begin{theorem}[C038]
With $J2 = \begin{smallmatrix}0&1\\-1&0\end{smallmatrix}$, $J2^2 = -I$, and
$\begin{smallmatrix}a&b\\-b&a\end{smallmatrix} = aI + b\,J2$:
(a) the symmetric part ($aI$) and the antisymmetric part ($bJ$) are orthogonal
in the Frobenius inner product --- orthogonality is a property of the $90^\circ$
rotation, not information independence; (b) the block acts on a real vector as
complex multiplication with conjugation:
$\begin{smallmatrix}a&b\\-b&a\end{smallmatrix}\cdot[x;y] =
[\mathrm{Re}((a-bi)(x+yi)), \mathrm{Im}((a-bi)(x+yi))]$; (c) the block is
determined by $(a,b)$ (bijection to $\mathbb C$): 4 stored entries carry
2 degrees of freedom --- 50\% square redundancy.
\end{theorem}

The SSM state $M$ (dimension $[n_v, d_k, d_k]$ real) is therefore a block
diagonal of complex numbers; the complex compactification
$[n_v, d_k/2, d_k/2]$ is lossless (R048 compactification). Empirically, the
complex-structure deviation ($c = -b$, $a = d$) is $\approx 3\times10^{-5}$
across all 64 layers, and the single-step output error of the compactified
forward pass is $2.4\times10^{-4}$.

\section{Bridging to inference}

The formalization is implemented and measured in an assembly inference backend
(the artifact includes the measured results):

\paragraph{Input chain is bit-exact.} The kernel chain
(per-block weight locking $\to$ window sliding with pre-locking $\to$
integer-MAC convolution $\to$ L2 compactification) matches the float reference
exactly on the input chain: window codes, convolution codes, and the L2-compactified
q/k codes all have zero difference; the single-layer output relative error is
$< 10^{-4}$ (correlation $> 0.9999$), and the layer output argmax is identical.

\paragraph{Locked-domain decisions are preserved.} End-to-end on real
Bonsai-27B Q2\_0 weights (SSM layers replaced by the assembly chain, all other
layers float), the locked domain (top1--top2 probability gap $> 0.7$) is
preserved on 6/6 generated tokens; the argmax is preserved on 2/6 tokens
(consistent with the basepoint-relative semantics: relative displacement is
invariant under basepoint change (C007), while absolute logits drift).

\paragraph{Q2\_0 grid in the wild.} The Bonsai-27B Q2\_0 weights indeed use
the grid $\{-1,0,1,2\}$ with per-block scales, matching C037's semantic: the
weights are basepoint-cluster symbols (C028), and their dequantization
$(\text{code}-1)\cdot d$ recovers the signed values.

\section{Related work and conclusion}

Basepoint-relative algebra formalized here builds on classical heap/torsor
theory \cite{bergman2009universal}; the operator-semantics layer (C026, C028--C031, C033,
C035--C036) formalizes classical facts (ring theory, unit-circle geometry,
dihedral groups, trigonometry) with a novel basepoint-relative reading; C037
is a direct mapping between the Q2\_0 quantization format and the four-phase
structure (novelty unassessed); C038 formalizes the classical complex-structure
block with the 50\%-redundancy observation and empirical validation.

The Lean formalization is developed in mathlib \cite{mathlib2026}; the
quantization format formalized in C037 is the GGUF Q2\_0 scheme \cite{gguf2026}.
The Riemann-direction formalization (C011--C025, previously submitted) and the
pat theorem chain (R047--R232) provide the background theory; this paper is
self-contained in its claims C026--C038.

\paragraph{Conclusion.} Operator semantics for quantized recurrent inference
can be formalized: every operator has its own basepoint; the Q2\_0 grid is a
two-sign-axis symbol encoding of the four phases; the SSM state is a complex
structure with 50\% redundancy. The formalization is machine-checked, and an
assembly implementation of the same semantics is bit-exact on the input chain
and preserves locked-domain decisions on real weights.

\section*{Theorem inventory (C026--C038)}

\begin{center}
\small
\begin{tabular}{lll}
\toprule
Claim & File & Key theorems \\
\midrule
C026 & BasepointMove.lean & add\_iter\_two, mul\_iter\_one\_eq\_add\_iter\_two,\\
     & & mul\_iter\_two, core, basepointMove, addOpInvariant,\\
     & & mulOpAffine, centerNotOrigin \\
C028 & OperatorFamily.lean & neg\_involution (operator basepoints) \\
C029 & HighDimProjection.lean & real\_proj\_range, real\_proj\_extreme\_pos/neg,\\
     & & real\_proj\_extremes, quad\_phase\_proj \\
C030 & SymplecticProjection.lean & J²=-I, phase sum/cancellation \\
C031 & Discretization.lean & quad\_phase\_card, recip\_is\_conjugate,\\
     & & conjugate\_anticommutes\_J, discrete\_projection\_cancel,\\
     & & basepoint\_cluster \\
C033 & DualProjection.lean & dual\_projection\_decomposition, J\_swaps\_projections,\\
     & & dual\_projection\_orthogonal, dual\_projection\_quad\_discrete \\
C035 & MatrixOpsPat0.lean & matrix operator theorems \\
C036 & TrigPat0.lean & cos\_even\_direction\_loss, sin\_odd\_direction\_keep,\\
     & & cos\_not\_injective, rot\_keeps\_norm, cos/sin\_basepoint,\\
     & & quad\_phase\_grid\_*, rotJ\_sq\_neg\_id, rotJ\_pow4\_id,\\
     & & rot90\_phase\_advance \\
C037 & Q2GridSymbolQuantization.lean & q2\_phi\_on\_phases, q2\_sign\_axes, q2\_iter\_axes,\\
     & & pat\_three\_is\_real\_projection, q2\_four\_vs\_pat\_three,\\
     & & four\_phase\_mul, four\_phase\_sum \\
C038 & ComplexStateCompactification.lean & j2\_sq, complex\_block\_decomp,\\
     & & sym\_skew\_orthogonal\_frobenius,\\
     & & complex\_block\_mul\_equiv, complex\_block\_bijective \\
\bottomrule
\end{tabular}
\end{center}

\bibliographystyle{ACM-Reference-Format}
\bibliography{refs}

\end{document}
